Let $\mathfrak{A}$ be a commutative ring with an identity. Consider the ring of functions $\mathfrak{A} = (\mathfrak{A}, \mathfrak{A})$. The identity function $s \mapsto s$ commutes with the constant functions. The elements of the subring $\mathfrak{A}[s]$ generated by the constant functions and the identity function are called **polynomial functions in one variable**. If $f(x) = a_0 + a_1x + \cdots + a_nx^n \in \mathfrak{A}[x]$ where $x$ is transcendental, then $f(s)$ maps $s$ to $a_0 + a_1s + \cdots + a_ns^n \in \mathfrak{A}$, and $\mathfrak{A}[s]$ is the totality of these functions.

More generally, polynomial functions in $r$ variables are the elements of the ring $\mathfrak{A}[s_1, s_2, \cdots, s_r]$ generated by the constant functions and the $r$ functions $s_i$. The $s_i$ commute with each other and with the constant functions. If $f(x_1, \cdots, x_r) \in \mathfrak{A}[x_1, \cdots, x_r]$ with $x_i$ algebraically independent, then $f(s_1, \cdots, s_r)$ is a polynomial function and every polynomial function is obtained in this way.

If $\mathfrak{A}$ is a finite ring of $q$ elements $a_j$, then
$$h(s_i) = (s_i - a_1)(s_i - a_2) \cdots (s_i - a_q) = 0,$$
so the functions $s_1, \cdots, s_r$ are algebraic relative to the subring of constant functions.
